\exercisegroup

\begin{enumerate}
\def\labelenumi{\arabic{enumi})}
\item
  Let G be the topological group O(2, R). We have \(\pi_{0}(G) \simeq Z/2Z\) and \(\pi_{1}(G) \simeq Z\) ; the nontrivial element of \(\pi_{0}(G)\) acts on \(\pi_{1}(G)\) by the mapping \(t \mapsto -t\) .
\item
  We denote by H the algebra of Hamilton quaternions (TG, VIII, p. 4) and \((1, i, j, k)\) its canonical basis. Let N denote the norm quadratic form of H; the set of quaternions of norm 1 is the sphere \(S_{3}\) . Let \(H_{0}\) be the real vector subspace of quaternions with zero trace and \(N_{0}\) the restriction to \(H_{0}\) of N.
\end{enumerate}

\begin{enumerate}
\def\labelenumi{\alph{enumi})}
\item
  Prove that the homomorphism in Proposition 4 of TG, VIII, p. 4 makes the sphere \(\mathbf{S}_3\) a covering of degree 2 of \(\mathbf{SO}(\mathrm{N}_0)\) .
\item
  Hence deduce that the Poincaré group of SO(3, R) is isomorphic to Z/2Z. Explicitly describe a loop in \(I_{3}\) in SO(3, R) whose strict homotopy class is the unique nontrivial element of this Poincaré group.
\item
  Let \(\varphi\) be the map from \(\mathbf{S}_3 \times \mathbf{S}_3\) into SO(4, R) given by \(\varphi(\mathbf{q}_1, \mathbf{q}_2)(\mathbf{x}) = \mathbf{q}_1 \mathbf{x} \mathbf{q}_2^{-1}\). Prove that \(\varphi\) makes \(\mathbf{S}_3 \times \mathbf{S}_3\) a covering of degree 2.
\item
  Deduce that the Poincaré group of SO(4, R) is isomorphic to Z/2Z.
\end{enumerate}

\begin{enumerate}
\def\labelenumi{\arabic{enumi})}
\setcounter{enumi}{2}
\tightlist
\item
  Let n be a natural number. Let the topological group \(\mathrm{SO}(n+1,\mathbf{R})\) act on the sphere \(S_{n}\) of \(R^{n+1}\) .
\end{enumerate}

\begin{enumerate}
\def\labelenumi{\alph{enumi})}
\item
  Prove that the fixator of the vector \(e_{n+1}\) is identified with the topological group \(\mathrm{SO}(n,\mathbf{R})\) .
\item
  Let \(f \colon \mathbf{I} \to \mathbf{S}_n\) be a continuous map. Assume that \(f(t) \neq (0, \ldots, 0, 1)\) for all \(t \in \mathbf{I} .\) Prove that there exists a continuous map \(g \colon \mathbf{I} \to \mathrm{SO}(n + 1, \mathbf{R})\) such that \(f(t) = g(t) \cdot \mathbf{e}_{n+1}\) for all \(t \in \mathbf{I} \).
\item
  If \(n \geqslant 3\) , prove that the inclusion of \(\mathrm{SO}(n, \mathbf{R})\) into \(\mathrm{SO}(n+1, \mathbf{R})\) induces an isomorphism between Poincaré groups.
\end{enumerate}

\begin{enumerate}
\def\labelenumi{\arabic{enumi})}
\setcounter{enumi}{3}
\tightlist
\item
  Let n be a natural number. Let \(B_{0}\) be the subgroup of \(\mathbf{SL}(n,\mathbf{R})\) consisting of upper triangular matrices whose diagonal coefficients are all strictly positive.
\end{enumerate}

\begin{enumerate}
\def\labelenumi{\alph{enumi})}
\item
  Prove that the map from \(\mathrm{SO}(n,\mathbf{R})\times\mathrm{B}_{0}\) into \(\mathbf{SL}(n,\mathbf{R})\) given by \((g,u)\mapsto g\cdot u\) is a homeomorphism (cf. INT, VII, p. 91, prop. 7).
\item
  Prove that the inclusion of SO(n, R) into SL(n, R) induces an isomorphism between Poincaré groups.
\end{enumerate}

\begin{enumerate}
\def\labelenumi{\arabic{enumi})}
\setcounter{enumi}{4}
\tightlist
\item
  Let G be a connected loopable topological group, let \((\widetilde{\mathrm{G}}, p)\) be a universal cover of G; we endow \(\widetilde{G}\) with the topological group structure for which p is a group homomorphism. We denote by Z the center of G, by e its identity element, \(\widetilde{Z}\) the center of \(\widetilde{G}\) and \(\widetilde{e}\) its identity element.
\end{enumerate}

\begin{enumerate}
\def\labelenumi{\alph{enumi})}
\item
  Let \(f\) be a continuous automorphism of the group G. Prove that there exists a unique continuous map \(\widetilde{f} \colon \widetilde{\mathrm{G}} \to \widetilde{\mathrm{G}}\) such that \(\widetilde{f}(\widetilde{e}) = \widetilde{e}\) and \(p \circ \widetilde{f} = f \circ p\). Prove that \(\widetilde{f}\) is a group automorphism. Prove that the map \(\varphi \colon f \mapsto \widetilde{f}\) is an injective group homomorphism from \(\operatorname{Aut}(\mathbf{G})\) into \(\operatorname{Aut}(\widetilde{\mathbf{G}})\) .
\item
  We denote \(\operatorname{Out}(\mathbf{G})\) the quotient of the group \(\operatorname{Aut}(\mathbf{G})\) by the subgroup of inner automorphisms; one defines \(\operatorname{Out}(\widetilde{\mathbf{G}})\) analogously. Prove that the homomorphism \(\varphi\), induced by passing to quotients, is an injective group homomorphism \(\psi\colon\operatorname{Out}(\mathbf{G})\to\operatorname{Out}(\widetilde{\mathbf{G}})\) .
\end{enumerate}

\begin{enumerate}
\def\labelenumi{\arabic{enumi})}
\setcounter{enumi}{5}
\tightlist
\item
  Let G and H be connected real Lie groups, and let \(f: G \to H\) be a surjective morphism of Lie groups. We denote by \(Z_{G}\) the center of G and by \(Z_{H}\) that of H.
\end{enumerate}

\begin{enumerate}
\def\labelenumi{\alph{enumi})}
\item
  Prove the equivalence of the following assertions: (i) The map f makes G a covering of H; (ii) \(\operatorname{Ker}(f)\) is discrete; (iii) The morphism \(\mathrm{L}(f)\colon\mathrm{L}(\mathrm{G})\to\mathrm{L}(\mathrm{H})\) induced from f by passing to Lie algebras is an isomorphism.
\item
  Prove that then \(f(\mathrm{Z}_{\mathrm{G}})=\mathrm{Z}_{\mathrm{H}}\) and that \(f^{-1}(\mathrm{Z}_{\mathrm{H}})=\mathrm{Z}_{\mathrm{G}}\) .
\end{enumerate}

\begin{enumerate}
\def\labelenumi{\arabic{enumi})}
\setcounter{enumi}{6}
\tightlist
\item
  Let G be a real Lie group, denote by \(G_{0}\) its neutral component, \(Z_{0}\) the center of \(G_{0}\) and \(\pi_{0}(G)=G/G_{0}\) . Let \((\widetilde{G}_{0},p_{0})\) be a universal covering of \(G_{0}\) , \(\widetilde{Z}_{0}\) the center of \(\widetilde{G}_{0}\) and \(\widetilde{e}\) its identity element.
\end{enumerate}

\begin{enumerate}
\def\labelenumi{\alph{enumi})}
\item
  Let \(\theta\colon\pi_{0}(\mathrm{G})\to\mathrm{Out}(\mathrm{G}_{0})\) be the homomorphism associated with the extension \(\mathcal{E}\colon\mathrm{G}_{0}\to\mathrm{G}\to\pi_{0}(\mathrm{G})\) (A, X, §7, exerc. 5). Prove that the cohomology class \(\omega(\pi_{0}(\mathrm{G}),\mathrm{Z}_{0},\theta)\) of \(\mathrm{H}^{3}(\pi_{0}(\mathrm{G}),\mathrm{Z}_{0})\) defined in loc. cit. is zero.
\item
  Assume that there exists a Lie group \(\widetilde{\mathrm{G}}\), a surjective map \(p\colon\widetilde{\mathrm{G}}\to\mathrm{G}\), and an isomorphism \(j\) of \(\widetilde{\mathrm{G}}_{0}\) onto the neutral component of \(\widetilde{\mathrm{G}}\) such that \(p_{0}=p\circ j\). Show that we have \(\omega(\pi_{0}(\mathrm{G}),\widetilde{\mathrm{Z}}_{0},\psi\circ\theta)=0\) in \(\mathrm{H}^{3}(\pi_{0}(\mathrm{G}),\widetilde{\mathrm{Z}}_{0})\) , where \(\psi\colon\mathrm{Out}(\mathrm{G}_{0})\to\mathrm{Out}(\widetilde{\mathrm{G}}_{0})\) is the homomorphism defined in exerc. 5.
\item
  Assume that \(\omega(\pi_{0}(\mathrm{G}),\widetilde{\mathrm{Z}}_{0},\psi\circ\theta)=0\) . Let \(\widetilde{E}\colon\widetilde{G}_{0}\to E\to\pi_{0}(G)\) be an extension of \(\pi_{0}(G)\) by \(\widetilde{G}_{0}\) whose associated homomorphism is equal to \(\psi\circ\theta\). Prove that \(\mathrm{E}/\pi_{1}(\mathrm{G}_{0})\) is an extension of \(\pi_{0}(G)\) by \(G_{0}\) whose associated homomorphism is equal to \(\theta\) .
\end{enumerate}

Now let \(\gamma\) be the unique element of \(\mathrm{H}^2 (\pi_0(\mathrm{G}),\mathrm{G}_0)\) such that \(\gamma \cdot [\mathrm{E} / \pi_1(\mathrm{G}_0)] = [\mathrm{G}]\) (loc. cit., e)). Consider the homomorphism \(\delta \colon \mathrm{H}^2 (\pi_0(\mathrm{G}),\mathrm{Z}_0)\to\) \(\mathrm{H}^3 (\pi_0(\mathrm{G}),\pi_1(\mathrm{G}))\) associated with the exact sequence of abelian groups \(1\to \pi_1(\mathrm{G}_0)\to\) \(\widetilde{\mathbf{Z}}_0\to \mathbf{Z}_0\to 1\) (exerc. 6). Show that the element \(\delta (\gamma)\) of \(\mathrm{H}^3 (\pi_0(\mathrm{G}),\pi_1(\mathrm{G}))\) does not depend on the choice of the extension \(\widetilde{\mathcal{E}}\) ; it is denoted \(\delta (\mathrm{G})\) .

d) Assume that \(\omega (\pi_0(\mathrm{G}),\widetilde{\mathrm{Z}}_0,\psi \circ \theta) = 0\) . For there to exist a Lie group \(\widetilde{\mathrm{G}}\), a surjective map \(p\colon \widetilde{\mathrm{G}}\to \mathrm{G}\), and an isomorphism \(j\) of \(\widetilde{\mathrm{G}}_0\) onto the neutral component of \(\widetilde{G}\) such that \(p_{0}=p\circ j\) , it is necessary and sufficient that \(\delta(\mathrm{G})=0.\) (5)

(5) For examples, see the article by R. L. TAYLOR, “Covering groups of nonconnected topological groups,” Proc. Amer. Math. Soc. 5 (1954), pp. 753–768.

\begin{enumerate}
\def\labelenumi{\arabic{enumi})}
\setcounter{enumi}{7}
\tightlist
\item
  Let G be a Lie group and let H be a closed subgroup of G; we denote by \(H_{0}\) the neutral component of H and \(\pi_{0}(H)\) the group \(H/H_{0}\) . Let us also denote by p the canonical surjection from G to G/H.
\end{enumerate}

\begin{enumerate}
\def\labelenumi{\alph{enumi})}
\item
  The map p has the path lifting property.
\item
  If H is connected, the homomorphism \(p_{*}\colon\pi_{1}(\mathrm{G},e)\to\pi_{1}(\mathrm{G}/\mathrm{H},p(e))\) is surjective.
\item
  There exists a unique map \(\varphi\colon\pi_{1}(\mathrm{G/H},p(e))\to\pi_{0}(\mathrm{H})\) which, for every path \(\widetilde{c}\) with source \(e\) in G whose term belongs to H, associates to the class of the loop \(p\circ\widetilde{c}\) in G/H the class of \(\widetilde{c}(1)\) modulo \(\mathrm{H}_{0}\) ; it is a group homomorphism.
\item
  If G is simply connected, then \(\varphi\) is an isomorphism. In particular, H is connected if and only if G/H is simply connected.
\end{enumerate}

\begin{enumerate}
\def\labelenumi{\arabic{enumi})}
\setcounter{enumi}{8}
\tightlist
\item
  Let G be a connected Lie group, let \((\widetilde{\mathbf{G}}, q)\) be a universal cover of G and let \(\widetilde{e}\) be the identity element of \(\widetilde{G}\). Let H be a closed connected subgroup of G; we denote by \(i: H \to G\) the canonical injection and \(p: G \to G/H\) the canonical surjection.
\end{enumerate}

\begin{enumerate}
\def\labelenumi{\alph{enumi})}
\item
  Let \(H_{1}\) be the neutral component of \(\overline{q}^{-1}(H)\) and let \(q_{1}\colon\widetilde{\mathrm{G}}/\mathrm{H}_{1}\to\mathrm{G}/\mathrm{H}\) be the map induced by q by passing to the quotients. Prove that \(q_{1}\) makes \(\widetilde{G}/H_{1}\) a covering of G/H whose fiber at \(p(e)\) is isomorphic to \(\pi_{1}(\mathrm{G}/\mathrm{H},p(e))\) .
\item
  Deduce that there exists an exact sequence of groups
\end{enumerate}

\[
\pi_ {1} (\mathrm{H} _ {1}, \widetilde {e}) \xrightarrow {q _ {*}} \pi_ {1} (\mathrm{H}, e) \xrightarrow {i _ {*}} \pi_ {1} (\mathrm{G}, e) \xrightarrow {p _ {*}} \pi_ {1} (\mathrm{G} / \mathrm{H}, p (e)).
\]

\begin{enumerate}
\def\labelenumi{\arabic{enumi})}
\setcounter{enumi}{9}
\tightlist
\item
  \begin{enumerate}
  \def\labelenumii{\alph{enumii})}
  \tightlist
  \item
    Prove that, for every integer \(n \geqslant 2\) , the group \(\mathrm{SU}(n, \mathbf{C})\) is simply connected. (The space \(\mathrm{SU}(2, \mathbf{C})\) is homeomorphic to \(S_{3}\). Let \(\mathrm{SU}(n, \mathbf{C})\) act on the unit sphere of \(C^{n}\), identify the stabilizer of the point \((0, \ldots, 0, 1)\) with the group \(\mathrm{SU}(n - 1, \mathbf{C})\), and use Exercise 9.)
  \end{enumerate}
\end{enumerate}

\begin{enumerate}
\def\labelenumi{\alph{enumi})}
\setcounter{enumi}{1}
\tightlist
\item
  Prove that for every integer \(n \geqslant 2\) , the group \(\mathbf{SL}(n, \mathbf{C})\) is simply connected. (Let \(B_{+}\) be the subgroup of \(\mathbf{SL}(n, \mathbf{C})\) consisting of upper triangular matrices whose diagonal coefficients are real numbers \textgreater{} 0. Show that the map from \(\mathrm{SU}(n, \mathbf{C}) \times \mathrm{B}_{+}\) into \(\mathbf{SL}(n, \mathbf{C})\) given by \((g, u) \mapsto g \cdot u\) is a homeomorphism.)
\end{enumerate}

